%! AP Statistics — Unit 1, Lesson 1.7 — Homework ANSWER KEY
\documentclass[10pt]{article}
\usepackage{apstats-article}
\usepackage{apstats-key}

%=============================================================================
\begin{document}

\pageheader{Unit 1, Lesson 1.7}{Homework --- Answer Key}

\vspace{0.10in}

\begin{tcolorbox}[colback=sky, colframe=navy, boxrule=0.8pt, arc=2mm,
    left=3mm, right=3mm, top=1.5mm, bottom=1.5mm]
\small For all problems: show calculations, interpret statistics in context,
and use complete sentences for AP-style questions.
\end{tcolorbox}

\vspace{0.12in}

% ════════════════════════════════════════════════════════════════════════════
% PART A
% ════════════════════════════════════════════════════════════════════════════
\boxguard
\begin{blurbbox}[Part A \quad Calculate Summary Statistics]{navylight}
\small

A survey asked 10 students: ``How many hours did you use your phone last week?''
The sorted responses are:
\[3,\quad 5,\quad 7,\quad 8,\quad 9,\quad 10,\quad 12,\quad 14,\quad 16,\quad 22\]

\vspace{0.10in}

\begin{enumerate}[leftmargin=1.5em, itemsep=8pt]

  \item Find the \textbf{mean}. Show your calculation.
\begin{work}
  \bar{x} &= \dfrac{3+5+7+8+9+10+12+14+16+22}{10} \\
          &= \dfrac{106}{10} \\
          &= 10.6 \text{ hours}
\end{work}

  \item Find the \textbf{median}. Which two values did you average?
\begin{work}
  n &= 10 \text{ (even), so average the 5th and 6th values} \\
  M &= \dfrac{9+10}{2} \\
    &= 9.5 \text{ hours}
\end{work}

  \item Find $Q_1$, $Q_3$, and the \textbf{IQR}.\\[5pt]
        $Q_1 = \ans{7}$ \quad $Q_3 = \ans{14}$ \quad IQR $= \ans{7}$

  \item Use the fence rule to check whether 22 is an outlier. Show both fences.\\[5pt]
        Lower fence: \ans{$7-1.5(7)=-3.5$} \quad Upper fence: \ans{$14+1.5(7)=24.5$}\\[4pt]
        Is 22 an outlier? \ans{No --- $22 < 24.5$}

  \item Is the mean or median a better description of the ``typical'' hours of
        phone use? Write one complete sentence with your reasoning.\\[5pt]
        \ansline{The \textbf{median (9.5 hours)} is the better description, because the}\\[1pt]
        \ansline{distribution is slightly skewed right (mean $10.6 >$ median $9.5$)}\\[1pt]
        \ansline{and the median resists the higher values in the tail.}

\end{enumerate}

\end{blurbbox}

\vspace{0.12in}

% ════════════════════════════════════════════════════════════════════════════
% PART B
% ════════════════════════════════════════════════════════════════════════════
\boxguard
\begin{blurbbox}[Part B \quad Identify an Outlier \& Choose Statistics]{greenacc}
\small

A class of 12 students took a test. The sorted scores are:
\[45,\quad 62,\quad 65,\quad 68,\quad 70,\quad 72,\quad 74,\quad 75,\quad 77,\quad 80,\quad 82,\quad 88\]

\begin{enumerate}[leftmargin=1.5em, itemsep=8pt]

  \item Find the median, $Q_1$, $Q_3$, and IQR. Show your work.
\begin{work}
  M   &= \dfrac{72+74}{2} = 73 \\
  Q_1 &= \dfrac{65+68}{2} = 66.5 \\
  Q_3 &= \dfrac{77+80}{2} = 78.5 \\
  \text{IQR} &= 78.5 - 66.5 = 12
\end{work}

  \item Use the fence rule to check whether 45 is an outlier.
\begin{work}
  \text{lower fence} &= Q_1 - 1.5(\text{IQR}) \\
                     &= 66.5 - 1.5(12) \\
                     &= 48.5
\end{work}
        Is 45 an outlier? \ans{Yes --- $45 < 48.5$}

  \item The mean of this dataset is 71.5 and the median is 73.
        What does mean $<$ median suggest about the shape?\\[5pt]
        \ansline{Mean $(71.5) <$ median $(73)$ suggests the distribution is}\\[1pt]
        \ansline{\textbf{skewed left}; the low outlier 45 pulls the mean down.}

  \item Which pair of statistics should be reported to describe this class's
        test performance? Name the pair and explain why.\\[5pt]
        \ansline{\textbf{Median + IQR}, because the distribution is skewed left and has}\\[1pt]
        \ansline{an outlier (45); both are resistant to that extreme score.}

\end{enumerate}

\end{blurbbox}

\vspace{0.12in}

% ════════════════════════════════════════════════════════════════════════════
% PART C — AP-STYLE
% ════════════════════════════════════════════════════════════════════════════
\boxguard
\begin{blurbbox}[Part C \quad AP-Style Free Response]{redacc}
\small

A dataset has mean 82 and median 91.

\begin{enumerate}[leftmargin=1.5em, itemsep=8pt]

  \item What does the relationship between the mean and median suggest about
        the shape of this distribution? Explain your reasoning in one or two
        complete sentences.\\[5pt]
        \ansline{Because the mean (82) is less than the median (91), the distribution}\\[1pt]
        \ansline{is likely \textbf{skewed left}: a few unusually low values pull the mean}\\[1pt]
        \ansline{below the center, while the median resists this effect.}

  \item Which measure of center --- mean or median --- is more appropriate for
        describing a typical value in this distribution? Justify your answer.\\[5pt]
        \ansline{The \textbf{median (91)}, because the distribution is skewed left; the}\\[1pt]
        \ansline{mean is dragged downward by the extreme low values.}

\end{enumerate}

\end{blurbbox}

\end{document}
